\section{积分的微分：从平均值恢复点值}\label{sec:04-01}

\subsection{Vitali 覆盖与 Hardy--Littlewood 极大函数}\label{subsec:4-1-1}
\index{Hardy--Littlewood 极大函数}\index{Vitali 覆盖}

设 $f$ 是局部可积函数。若只看一个固定半径 $r$，球平均
\[
  \frac{1}{m(B(x,r))}\int_{B(x,r)}|f(y)|\,\dd y
\]
并不难处理。困难来自“同时看所有半径”：一个点只要在某个尺度上平均过大，就应当算作坏点。
每个坏点都带来一个见证球，而这些球可能层层重叠。把它们的体积直接相加，会把同一份函数质量
重复计算任意多次。我们需要先从见证球中挑出互不相交的一族，同时又不丢掉原来的坏点。

\begin{definition}[中心极大函数]\label{def:4-1-1-1}
对 $f\in L^1_{\mathrm{loc}}(\R^n)$，定义其\emph{中心 Hardy--Littlewood 极大函数}
\[
  Mf(x)=\sup_{r>0}\frac{1}{m(B(x,r))}
                \int_{B(x,r)}|f(y)|\,\dd y .
\]
这里及下文 $B(x,r)$ 表示 Euclidean 开球，并记 $\omega_n=m(B(0,1))$。
\end{definition}

这个定义中的上确界确实给出可测函数。固定 $r>0$，令
$I_r(x)=\int_{B(x,r)}|f|$。当 $x_k\to x$ 时，可把所有相关球限制在某个有界集合 $Q$ 内；
$|f|\1_Q\in L^1$。若 $h_k=|x_k-x|<r$，则
\[
 B(x_k,r)\mathbin\triangle B(x,r)
 \subset B(x,r+h_k)\setminus\overline{B(x,r-h_k)},
\]
因而
\[
 m\bigl(B(x_k,r)\mathbin\triangle B(x,r)\bigr)
 \le \omega_n\bigl((r+h_k)^n-(r-h_k)^n\bigr)\longrightarrow0.
\]
由积分的绝对连续性，
\[
 |I_r(x_k)-I_r(x)|
 \le \int_{B(x_k,r)\triangle B(x,r)}|f(y)|\,\dd y\longrightarrow0.
\]
因此每个固定半径的平均关于 $x$ 连续，连续函数族的上确界 $Mf$ 下半连续；特别地，
$\{Mf>\lambda\}$ 是开集。另一方面，若 $r_k\to r>0$，则
\[
 B(x,r_k)\triangle B(x,r)
 \subset B\bigl(x,r+|r_k-r|\bigr)
       \setminus\overline{B\bigl(x,r-|r_k-r|\bigr)},
\]
所以同一估计证明球积分关于 $r$ 连续。因此定义中的上确界也可只取
$r\in\Q_{>0}$。

\begin{definition}[Vitali 球族]\label{def:4-1-1-2}
称球族 $\mathcal V$ 以 Vitali 意义覆盖 $E\subset\R^n$，若对每个 $x\in E$ 和
$\delta>0$，$\mathcal V$ 中都有一个以 $x$ 为中心、半径小于 $\delta$ 的球。
\end{definition}

例如，开集 $G$ 内所有满足 $B(x,r)\Subset G$ 的中心球构成 $G$ 的 Vitali 覆盖。
相反，只允许一个固定半径的中心球族不满足“半径任意小”的要求。
下面的选择引理只要求球的半径有统一上界，并不要求半径可以任意小。任意小半径条件属于微分基的
局部性质，二者应予区分。

\begin{definition}[弱 $L^1$ 空间]\label{def:4-1-1-3}
可测函数 $g$ 属于弱 $L^1$，记作 $g\in L^{1,\infty}$，若
\[
  \|g\|_{L^{1,\infty}}
  :=\sup_{\lambda>0}\lambda\,m\{x:|g(x)|>\lambda\}<\infty.
\]
这个泛函通常称为弱 $L^1$ 准范数；它一般不满足范数的三角不等式。
\end{definition}

先看一个能直接算出的极大函数。

\begin{example}[球的示性函数]\label{ex:4-1-1-1}
令 $f=\1_{B(0,1)}$。若 $x\in B(0,1)$，任意充分小的 $B(x,r)$ 都包含在单位球内，故
$Mf(x)=1$。若 $|x|>2$，取半径 $|x|+1$ 的中心球，它包含 $B(0,1)$，于是
\[
  Mf(x)\ge \frac{m(B(0,1))}{m(B(x,|x|+1))}=(|x|+1)^{-n}.
\]
反之，任何与 $B(0,1)$ 相交的 $B(x,r)$ 都满足 $r>|x|-1$，从而
\[
 \frac{m(B(x,r)\cap B(0,1))}{m(B(x,r))}
 \le (|x|-1)^{-n}.
\]
不相交的球给出平均值 $0$，所以
\[
 (|x|+1)^{-n}\le Mf(x)\le(|x|-1)^{-n}\qquad(|x|>2).
\]
特别地，$Mf(x)$ 在无穷远与 $|x|^{-n}$ 同阶；用极坐标看，
$\int_{|x|>2}|x|^{-n}\,\dd x$ 含有 $\int_2^\infty r^{-1}\,\dd r$，因而发散。
所以极大算子不可能把 $L^1$ 有界地映入 $L^1$。
\end{example}

\begin{example}[高度重叠的球族]\label{ex:4-1-1-2}
固定 $N\ge1$，取
\[
 B_k=B(0,2^{-k}),\qquad 1\le k\le N,
 \qquad f_N=\1_{B(0,2^{-N})}.
\]
这些球全部含有 $B_N$，因而
\[
 \sum_{k=1}^N\int_{B_k}f_N\,\dd x
 =N\,m(B_N),
 \qquad
 \int_{\bigcup_{k=1}^NB_k}f_N\,\dd x=m(B_N).
\]
直接对见证球求和已把同一份质量重复计算 $N$ 次。甚至把每个球复制有限多次也不改变
它们的并，却能任意增大体积之和。覆盖估计若不先做选择，便不可能只依赖被覆盖集合。
\end{example}

\begin{lemma}[$5r$ 覆盖引理]\label{lemma:4-1-1-1}
设 $\mathcal V$ 是 $\R^n$ 中半径有上界的球族。存在可数个两两不交的球
$B_j\in\mathcal V$，使得
\[
  \bigcup_{B\in\mathcal V}B\subset\bigcup_{j=1}^{\infty}5B_j.
\]
而且每个 $B\in\mathcal V$ 都与某个满足
$\operatorname{rad}(B_j)\ge\tfrac12\operatorname{rad}(B)$ 的 $B_j$ 相交。
\end{lemma}

\begin{proof}
先设 $\mathcal V$ 有限。按半径不增的次序检查各球：若当前球与已经选出的球都不交，就选入；
否则舍去。每个舍去的球 $B=B(x,r)$ 都与一个较早选出的
$B_j=B(x_j,r_j)$ 相交，并有 $r_j\ge r$。

一般情形取所有半径的一个上界 $R$，并按
\[
 \mathcal V_k=\{B\in\mathcal V:2^{-k-1}R<\operatorname{rad}(B)\le2^{-k}R\},
 \qquad k=0,1,\ldots,
\]
分层。从 $k=0$ 起，依次在 $\mathcal V_k$ 中选取一个极大的两两不交子族，同时要求它与此前层中
已选球都不交。这样的极大子族可由集合包含关系下的极大原理取得。任一两两不交的正半径球族都是
可数的：给每个球选一个不同的有理坐标点即可。因此全部选出的球可编号为 $(B_j)$。

若 $B\in\mathcal V_k$ 未被选入，极大性说明它与第 $k$ 层或更早一层的某个已选球 $B_j$ 相交。
若 $r,r_j$ 为二者半径，则 $r_j>2^{-k-1}R\ge r/2$。取交点 $w$。对任意 $z\in B$，
\[
 |z-x_j|\le |z-x|+|x-w|+|w-x_j|<2r+r_j\le5r_j,
\]
故 $B\subset5B_j$。已选球自身当然也包含在相应的 $5B_j$ 中，结论成立。
\end{proof}

有限球族的贪心选择甚至给出更好的放大常数；这里保留 $5$，是为了让有限与一般球族共用同一陈述。

\begin{theorem}[Hardy--Littlewood 弱型 $(1,1)$]\label{theorem:4-1-1-2}
对每个 $f\in L^1(\R^n)$ 和 $\lambda>0$，有
\[
 m\{Mf>\lambda\}\le \frac{5^n}{\lambda}\|f\|_1.
\]
特别地，$M:L^1(\R^n)\to L^{1,\infty}(\R^n)$。
\end{theorem}

\begin{proof}
记 $E_\lambda=\{Mf>\lambda\}$。上文已经证明它是开集。先取紧集
$K\subset E_\lambda$。对每个 $x\in K$，按极大函数的定义选 $r_x>0$，使
\[
  \int_{B(x,r_x)}|f|>\lambda m(B(x,r_x)).
\]
这些见证球覆盖 $K$，紧性给出有限子覆盖。对这有限族应用
\cref{lemma:4-1-1-1}，取得两两不交的 $B_1,\ldots,B_N$，且
$K\subset\bigcup_{j=1}^N5B_j$。于是
\[
 \begin{aligned}
 m(K)&\le\sum_{j=1}^N m(5B_j)
      =5^n\sum_{j=1}^Nm(B_j)\\
 &<\frac{5^n}{\lambda}\sum_{j=1}^N\int_{B_j}|f|
 \le\frac{5^n}{\lambda}\|f\|_1,
 \end{aligned}
\]
最后一步正是选择不交球的用途。Lebesgue 测度对开集内正则，故对所有紧集
$K\subset E_\lambda$ 取上确界即得所述估计。
\end{proof}

\begin{remark}[一维覆盖常数的改进]\label{ex:4-1-1-3}
在 $\R$ 上，球就是区间。若改用 rising sun lemma，可按开集分量选择区间，得到比 $5$ 更好的常数。
球覆盖证明则可直接用于随后的 $\R^n$ 微分定理；一维改进见本小节习题。
\end{remark}

\begin{proposition}[强型 $(p,p)$（选读）]\label{proposition:4-1-1-3}
若 $1<p\le\infty$，则存在只依赖 $n,p$ 的常数 $C_{n,p}$，使
\[
  \|Mf\|_p\le C_{n,p}\|f\|_p\qquad(f\in L^p(\R^n)).
\]
\end{proposition}

\begin{proof}
$p=\infty$ 时，每个球平均不超过 $\|f\|_\infty$。设 $1<p<\infty$。对固定
$\lambda>0$ 分解
 \[
 f=f_\lambda^{(0)}+f_\lambda^{(1)},\qquad
 f_\lambda^{(0)}=f\1_{\{|f|\le\lambda/2\}},\qquad
 f_\lambda^{(1)}=f\1_{\{|f|>\lambda/2\}}.
 \]
由于
\[
 \int_{\{|f|>\lambda/2\}}|f|
 \le (2/\lambda)^{p-1}\|f\|_p^p,
\]
$f_\lambda^{(1)}\in L^1$，所以可以对它使用弱型估计。
由 $M f_\lambda^{(0)}\le\lambda/2$ 及 $M(f+g)\le Mf+Mg$，
\[
 \{Mf>\lambda\}\subset
 \{M f_\lambda^{(1)}>\lambda/2\}.
\]
弱型估计给出
\[
 m\{Mf>\lambda\}
 \le\frac{2\cdot5^n}{\lambda}
       \int_{\{|f|>\lambda/2\}}|f(x)|\,\dd x.
\]
层蛋糕公式和 Tonelli 定理于是给出
\[
 \begin{aligned}
 \|Mf\|_p^p
 &=p\int_0^\infty\lambda^{p-1}m\{Mf>\lambda\}\,\dd\lambda\\
 &\le2\cdot5^np\int_{\R^n}|f(x)|
     \left(\int_0^{2|f(x)|}\lambda^{p-2}\,\dd\lambda\right)\dd x\\
 &=\frac{2^p5^np}{p-1}\|f\|_p^p.
 \end{aligned}
\]
取 $p$ 次方根即得结论。
\end{proof}

\begin{proposition}[局部极大函数]\label{proposition:4-1-1-4}
设 $K\subset\R^n$ 有界可测，$r_0>0$，并记
$K_{r_0}=\{y:\dist(y,K)<r_0\}$。若 $f\in L^1(K_{r_0})$，则
对 $x\in K$ 定义
\[
 M_{r_0}f(x)=\sup_{0<r<r_0}\frac1{m(B(x,r))}\int_{B(x,r)}|f(y)|\,\dd y
\]
满足
\[
 m\{x\in K:M_{r_0}f(x)>\lambda\}
 \le\frac{5^n}{\lambda}\int_{K_{r_0}}|f|.
\]
\end{proposition}

\begin{proof}
把 $f$ 在 $K_{r_0}$ 外延为零，所得函数记为 $F\in L^1(\R^n)$。若 $x\in K$ 且
$0<r<r_0$，则 $B(x,r)\subset K_{r_0}$，所以相应的 $f$ 平均等于 $F$ 平均。于是
$M_{r_0}f(x)\le MF(x)$，结论由全局弱型估计得到。
\end{proof}

若 $K\Subset\Omega$，取 $r_0<\dist(K,\Omega^c)$ 就得到完全在 $\Omega$ 内取平均的版本。
这与先把 $f$ 在 $\Omega$ 外延为零、再允许任意半径的算子不同：后者会把边界外的零也纳入平均。

\begin{figure}[htbp]
\centering
\begin{tikzpicture}[scale=.82]
  \foreach \x/\y/\r in {0/0/1.15,.7/.15/.78,-.65/.25/.63,.15/.7/.52,
                           1.35/.15/.46,-1.2/-.2/.40,.55/-.65/.34}
    \draw[BookInk!30] (\x,\y) circle (\r);
  \draw[BookAccent,very thick] (-.65,.25) circle (.63);
  \draw[BookAccent,very thick] (1.35,.15) circle (.46);
  \draw[BookAccent,very thick] (.55,-.65) circle (.34);
  \begin{scope}[xshift=6.1cm]
    \draw[BookAccent!55,thick,dashed] (-.65,.25) circle (3.15);
    \draw[BookAccent!55,thick,dashed] (1.35,.15) circle (2.30);
    \draw[BookAccent!55,thick,dashed] (.55,-.65) circle (1.70);
    \fill[BookAccent] (-.65,.25) circle (2pt);
    \fill[BookAccent] (1.35,.15) circle (2pt);
    \fill[BookAccent] (.55,-.65) circle (2pt);
  \end{scope}
  \draw[book arrow] (2.1,0)--(3.8,0)
    node[midway,above,book label]{选择不交球并放大};
\end{tikzpicture}
\caption{左侧见证球高度重叠；选择出的三个实线球互不相交。右侧的五倍球覆盖原球族的并，
而积分质量只在实线球上累计一次。}
\label{fig:4-1-vitali-5r}
\end{figure}

弱型不等式的常数 $5^n$ 来自球放大后的体积比，而不是某个隐藏的解析恒等式。它和上面使用的
示性函数点态估计都把“大于阈值”的集合转化为总质量除以阈值；不过前者还要先解决所有尺度见证球的重叠。
卷积近似、Poisson 平均以及调和函数的边界值问题也会遇到同时控制所有尺度的困难。概率论中的最大
不等式具有相似的尾界形式，但其对象是一族随机变量，而不是一族 Euclidean 球平均。

\begin{exercise}
证明有限球族按半径递减贪心选择后，其余每个球都包含在某个已选球的三倍球中。
\end{exercise}
\begin{exercise}
补全 $I_r(x)$ 连续性的证明，并由此证明 $Mf$ 下半连续。
\end{exercise}
\begin{exercise}
用 \cref{ex:4-1-1-1} 证明不存在常数 $C$ 使 $\|Mf\|_1\le C\|f\|_1$ 对所有
$f\in L^1(\R^n)$ 成立。
\end{exercise}
\begin{exercise}
重新计算 $M\1_{B(0,1)}$：分别用“包住单位球”的球与“只要相交便有半径下界”证明
\[
 (|x|+1)^{-n}\le M\1_{B(0,1)}(x)\le(|x|-1)^{-n}
 \qquad(|x|>2).
\]
\end{exercise}
\begin{exercise}
证明一维 rising sun lemma 的下列形式。若 $F\in C([a,b])$，则
\[
 G=\{x\in(a,b):F(x)<\max_{x\le y\le b}F(y)\}
\]
是可数个不交开区间之并；对每个不碰 $b$ 的分量 $(\alpha,\beta)$，有
$F(\alpha)=F(\beta)$。把它用于
$F(x)=\int_a^x|f(t)|\dd t-\lambda x$，推导一侧区间极大函数的弱型估计，并比较球覆盖证明中的常数。
\end{exercise}

极大函数已经把“某个半径上的平均异常”压缩成一个可测坏集，并控制了它的大小。下一步是先对连续函数
验证小球平均确实回到点值，再让稠密逼近和弱型估计共同处理一般局部可积函数。

\subsection{Lebesgue 微分定理与密度点}\label{subsec:4-1-2}
\index{Lebesgue 微分定理}\index{Lebesgue 点}\index{密度点}

积分会在固定尺度上平滑局部细节；Lebesgue 微分定理说明，当平均半径趋于零时，原函数仍可在几乎
每一点恢复。连续点上的结论直接来自连续性。对一般 $L^1_{\mathrm{loc}}$ 函数，则需要把连续函数的
计算通过稠密逼近与极大估计传递过去。

\begin{definition}[Lebesgue 点]\label{def:4-1-2-1}
设 $f\in L^1_{\mathrm{loc}}(\R^n)$ 是一个选定的可测代表。若 $f(x)$ 有限且
\[
 \lim_{r\downarrow0}\frac1{m(B(x,r))}
       \int_{B(x,r)}|f(y)-f(x)|\,\dd y=0,
\]
则称 $x$ 为 $f$ 的一个 \emph{Lebesgue 点}。
\end{definition}

\begin{definition}[密度点]\label{def:4-1-2-2}
设 $E\subset\R^n$ 可测。若
\[
 \lim_{r\downarrow0}\frac{m(E\cap B(x,r))}{m(B(x,r))}=1,
\]
则称 $x$ 为 $E$ 的密度点。若同一极限为 $0$，称 $x$ 为 $E$ 的密度零点。
\end{definition}

\begin{definition}[平均算子]\label{def:4-1-2-3}
对 $r>0$，定义
\[
 A_rf(x)=\frac1{m(B(0,r))}\int_{B(x,r)}f(y)\,\dd y.
\]
积分只依赖 $f$ 的局部 $L^1$ 等价类。
\end{definition}

\begin{example}[连续点]\label{ex:4-1-2-1}
若 $f$ 在 $x$ 连续，则给定 $\varepsilon>0$，存在 $\delta>0$，使得只要
$|y-x|<\delta$，就有 $|f(y)-f(x)|<\varepsilon$。因此 $0<r<\delta$ 时
\[
 \frac1{m(B(x,r))}\int_{B(x,r)}|f(y)-f(x)|\,\dd y\le\varepsilon.
\]
所以每个连续点都是 Lebesgue 点。这是后面逼近论证中负责提供正确点值的部分。
\end{example}

\begin{theorem}[Lebesgue 微分定理]\label{theorem:4-1-2-1}
若 $f\in L^1_{\mathrm{loc}}(\R^n)$，则 $f$ 的几乎每个点都是 Lebesgue 点。特别地，
\[
  A_rf(x)\longrightarrow f(x)\qquad\text{对几乎每个 }x
\]
成立。同样的结论对以 $x$ 为中心、边平行于坐标轴的立方体平均成立。
\end{theorem}

\begin{proof}
对 $f$ 的无限值点重新赋值为零不会改变其局部 $L^1$ 等价类，因此不妨设 $f$ 处处有限。固定正整数
$R$，令
\[
 K=B(0,R),\qquad F=f\1_{B(0,R+1)}\in L^1(\R^n).
\]
当 $x\in K$ 且 $0<r<1$ 时，$B(x,r)\subset B(0,R+1)$，故在这个球上 $F=f$。
定义
\[
 \Omega_f(x)=\limsup_{r\downarrow0}\frac1{m(B(x,r))}
             \int_{B(x,r)}|f(y)-f(x)|\,\dd y.
\]
固定 $x$ 时，球积分关于 $r$ 连续，所以可在趋于零的正有理半径上取上极限；由此
$\Omega_f$ 可测。这里对每个固定的有理 $r$，函数
$(x,y)\mapsto\1_{\{|y-x|<r\}}|f(y)-f(x)|$ 可测，参数积分的可测性保证相应球平均关于 $x$ 可测。

由 \cref{theorem:3-3-3-1} 中 $C_c(\R^n)$ 在 $L^1(\R^n)$ 中的稠密性，给定 $\eta>0$ 可取
$g\in C_c(\R^n)$ 使 $\|F-g\|_1<\eta$。对 $x\in K$ 及 $r<1$，三角不等式给出
\[
 \begin{aligned}
 \frac1{m(B(x,r))}\int_{B(x,r)}|f(y)-f(x)|\,\dd y
 &\le \frac1{m(B(x,r))}\int_{B(x,r)}|F(y)-g(y)|\,\dd y\\
 &\quad+\frac1{m(B(x,r))}\int_{B(x,r)}|g(y)-g(x)|\,\dd y\\
 &\quad+|g(x)-F(x)|.
 \end{aligned}
\]
中间一项因 $g$ 在 $x$ 连续而随 $r\downarrow0$ 趋于零，故
\[
  \Omega_f(x)\le M(F-g)(x)+|F(x)-g(x)|\qquad(x\in K).
\]
于是对 $\lambda>0$，由 \cref{theorem:4-1-1-2}，以及点态估计
\[
  \1_{\{|F-g|>\lambda/2\}}\le \frac{2}{\lambda}|F-g|,
\]
\[
 \begin{aligned}
 m\{x\in K:\Omega_f(x)>\lambda\}
 &\le m\{M(F-g)>\lambda/2\}
      +m\{|F-g|>\lambda/2\}\\
 &\le \frac{2(5^n+1)}{\lambda}\|F-g\|_1
 <\frac{2(5^n+1)}{\lambda}\eta.
 \end{aligned}
\]
$\eta$ 可任意小，故这个集合测度为零。对 $\lambda=1/k$ 取可数并，得到
$\Omega_f=0$ 几乎处处于 $K$。再令 $R$ 遍历正整数，得到几乎每点都是 Lebesgue 点。

最后，
\[
 |A_rf(x)-f(x)|\le\frac1{m(B(x,r))}
       \int_{B(x,r)}|f(y)-f(x)|\,\dd y,
\]
所以普通平均收敛是 Lebesgue 点结论的推论。对中心立方体 $Q(x,r)=x+[-r,r]^n$，有
$Q(x,r)\subset B(x,\sqrt n\,r)$，并且
$m(B(x,\sqrt n\,r))/m(Q(x,r))=\omega_n n^{n/2}/2^n$。因此
\[
 \frac1{m(Q(x,r))}\int_{Q(x,r)}|f(y)-f(x)|\,\dd y
 \le \frac{\omega_n n^{n/2}}{2^n}
 \frac1{m(B(x,\sqrt n r))}
 \int_{B(x,\sqrt n r)}|f(y)-f(x)|\,\dd y,
\]
右端趋于零，故立方体版结论成立。
\end{proof}

上述证明中，连续函数确定候选极限，$C_c$ 稠密性提供逼近，而极大弱型估计统一控制所有小半径上的
逼近误差。稠密性本身并不足以推出逐点结论。

\begin{corollary}[Lebesgue 密度定理]\label{corollary:4-1-2-2}
若 $E\subset\R^n$ 可测，则 $E$ 中几乎每点是 $E$ 的密度点，$E^c$ 中几乎每点是
$E$ 的密度零点。
\end{corollary}

\begin{proof}
$\1_E\in L^1_{\mathrm{loc}}$，并且
\[
 A_r\1_E(x)=\frac{m(E\cap B(x,r))}{m(B(x,r))}.
\]
对 $\1_E$ 应用 \cref{theorem:4-1-2-1}。在 $E$ 上其值为 $1$，在 $E^c$ 上其值为 $0$，
即得结论。
\end{proof}

\begin{example}[区间端点]\label{ex:4-1-2-2}
对 $E=(0,1)\subset\R$，若 $x\in(0,1)$，充分小的 $(x-r,x+r)$ 包含在 $E$ 中，密度为 $1$；
若 $x\notin[0,1]$，密度为 $0$。在 $x=0$，当 $0<r<1$ 时
\[
 \frac{m((0,1)\cap(-r,r))}{2r}=\frac12,
\]
$x=1$ 同样如此。两个端点正是零测的异常集。
\end{example}

\begin{example}[肥 Cantor 集]\label{ex:4-1-2-3}
从 $F_0=[0,1]$ 开始，在第 $k$ 步从 $F_{k-1}$ 的每个区间分量中央删去一个长度为 $4^{-k}$
的开区间，并令 $F=\bigcap_{k\ge1}F_k$。第 $k$ 步删去 $2^{k-1}$ 个区间，故删去的总长度为
\[
 \sum_{k=1}^\infty 2^{k-1}4^{-k}=\frac12,
\]
从而 $m(F)=1/2$。第 $k$ 级的 $2^k$ 个分量等长，其共同长度为
\[
 \ell_k=2^{-k}\left(1-\sum_{j=1}^k2^{j-1}4^{-j}\right)
 \le2^{-k}\longrightarrow0.
\]
若某个非退化开区间
$I$ 包含在 $F$ 中，它便包含在每个 $F_k$ 的某个分量内，这与分量长度趋于零矛盾。所以 $F$ 没有内点。
它又是闭集，因此每个 $x\in F$ 都是拓扑边界点。可是
\cref{corollary:4-1-2-2} 说明，$F$ 中几乎每个 $x$ 满足
\[
 \frac{m(F\cap(x-r,x+r))}{2r}\longrightarrow1.
\]
拓扑问题问“每个邻域是否碰到补集”，密度问题问“补集占邻域的多少”；这个例子使二者不能混为一谈。
\end{example}

\begin{proposition}[精确代表的约定]\label{proposition:4-1-2-4}
设 $f\in L^1_{\mathrm{loc}}(\R^n)$。在有限极限存在处令
\[
 \widetilde f(x)=\lim_{r\downarrow0}A_rf(x),
\]
在其他点令 $\widetilde f(x)=0$。则 $\widetilde f=f$ 几乎处处；而在极限存在处，
$\widetilde f(x)$ 只由 $f$ 的 $L^1_{\mathrm{loc}}$ 等价类决定。
\end{proposition}

\begin{proof}
若 $f=g$ 几乎处处，则对每个球 $B(x,r)$，二者积分相等，故 $A_rf(x)=A_rg(x)$；所有存在的
平均极限也相等。由 \cref{theorem:4-1-2-1}，对几乎每个 $x$，平均极限存在且等于所选代表
$f(x)$，所以 $\widetilde f=f$ 几乎处处。异常集上取零只是约定；积分数据并未在那里选出唯一点值。
\end{proof}

\begin{figure}[htbp]
\centering
\begin{tikzpicture}[x=1cm,y=1cm]
  \draw[->,BookInk] (-.2,0)--(8.1,0) node[right]{$x$};
  \draw[->,BookInk] (0,-.2)--(0,3.1) node[above]{$y$};
  \draw[BookInk!65,thick,smooth] plot coordinates
   {(0,1.2) (.6,1.5) (1.2,.8) (1.8,2.5) (2.4,.9) (3,1.8) (3.6,.6)
    (4.2,2.6) (4.8,1.0) (5.4,2.1) (6,1.2) (6.6,1.8) (7.4,1.4)};
  \draw[BookAccent,very thick,smooth] plot coordinates
   {(0,1.25) (.8,1.3) (1.6,1.35) (2.4,1.4) (3.2,1.45) (4,1.5)
    (4.8,1.55) (5.6,1.6) (6.4,1.65) (7.4,1.7)};
  \fill[BookWarm!18] (3.45,0) rectangle (4.55,2.85);
  \draw[BookWarm,decorate,decoration={brace,mirror,amplitude=4pt}]
    (3.45,-.05)--(4.55,-.05) node[midway,below=5pt]{$B(x,r)$};
  \node[BookAccent] at (6.7,2.45) {$g$：连续近似};
  \node[BookInk!75] at (6.8,.65) {$f$};
\end{tikzpicture}
\caption{连续近似 $g$ 的平均振荡随半径缩小而消失；$f-g$ 在所有半径上的误差由
$M(f-g)$ 统一控制。}
\label{fig:4-1-differentiation-approximation}
\end{figure}

微分理论还给出 Riemann 可积性的精确判据：决定因素不是不连续点是否有限或稠密，而是其集合是否
为零测集。

对有界 $f:[a,b]\to\R$ 定义其在 $x$ 的振荡
\[
 \omega_f(x)=\inf_{\delta>0}
 \sup\{|f(y)-f(z)|:y,z\in[a,b],\ |y-x|<\delta,\ |z-x|<\delta\}.
\]
$f$ 在 $x$ 连续当且仅当 $\omega_f(x)=0$。此外
$D_k=\{x:\omega_f(x)\ge1/k\}$ 是闭集：若 $\omega_f(x)<1/k$，存在一个邻域，使其中任一点
都继承同一个更大的振荡控制，故该邻域不碰 $D_k$。

\begin{theorem}[Riemann 可积判据]\label{theorem:4-1-2-3}
有界函数 $f:[a,b]\to\R$ Riemann 可积，当且仅当它的不连续点集合具有 Lebesgue 测度零。
\end{theorem}

\begin{proof}
先设 $f$ Riemann 可积。固定 $k\ge1$ 与 $\varepsilon>0$，取分割
$P:a=x_0<\cdots<x_N=b$ 使
\[
 U(f,P)-L(f,P)=\sum_{j=1}^N\operatorname{osc}(f;[x_{j-1},x_j])
                       (x_j-x_{j-1})<\frac{\varepsilon}{2k},
\]
其中 $\operatorname{osc}(f;I)=\sup_I f-\inf_I f$。若
$x\in D_k$ 且 $x$ 不是分割点，则含 $x$ 的分割小区间的振荡至少为 $1/k$。所有这类小区间的
总长度因此不超过
\[
 k\bigl(U(f,P)-L(f,P)\bigr)<\varepsilon/2.
\]
详细地，若 $x\in(x_{j-1},x_j)$，取
$0<\delta<\min\{x-x_{j-1},x_j-x\}$。由 $\omega_f(x)\ge1/k$，
\[
 \operatorname{osc}(f;[x_{j-1},x_j])
 \ge \sup_{\substack{|y-x|<\delta\\|z-x|<\delta}}|f(y)-f(z)|
 \ge\frac1k.
\]
因而把不含分割点的 $D_k$ 覆盖起来的小区间总长度确实至多为
$k(U(f,P)-L(f,P))$。
有限个分割点可由总长度小于 $\varepsilon/2$ 的有限开区间覆盖。所以 $D_k$ 的外测度小于
$\varepsilon$，即 $m(D_k)=0$。不连续点集为 $\bigcup_{k\ge1}D_k$，故它是零测集。

反之，设不连续点集 $D$ 测度为零，并令 $B=\sup_{[a,b]}|f|$。若 $B=0$，结论成立；以下设
$B>0$。给定 $\varepsilon>0$，由 Lebesgue 外正则性取开集 $G\supset D$，使
$m(G)<\varepsilon/(4B)$，并令 $K=[a,b]\setminus G$。对每个 $x\in K$，$f$ 在 $x$ 连续，
故可取 $\delta_x>0$，使当 $|y-x|<2\delta_x$ 时
\[
 |f(y)-f(x)|<\frac{\varepsilon}{4(b-a)}.
\]
有限个 $B(x_i,\delta_{x_i})$ 覆盖紧集 $K$。取分割 $P$ 的网格小于
$\min_i\delta_{x_i}$。若分割小区间 $I$ 与 $K$ 相交，取
$y\in I\cap K$，再选 $i$ 使 $y\in B(x_i,\delta_{x_i})$。对每个 $z\in I$，
$|z-x_i|\le |z-y|+|y-x_i|<2\delta_{x_i}$，故 $I\subset B(x_i,2\delta_{x_i})$。
于是这类 $I$ 上
\[
 \operatorname{osc}(f;I)<\frac{\varepsilon}{2(b-a)}.
\]
其余小区间不碰 $K$，因而包含于 $G$，它们的总长度至多 $m(G)$。所以
\[
 \begin{aligned}
 U(f,P)-L(f,P)
 &\le \frac{\varepsilon}{2(b-a)}(b-a)+2B\,m(G)\\
 &<\varepsilon.
 \end{aligned}
\]
这里第一类小区间的长度和至多 $b-a$；第二类小区间两两仅在端点相交，且其并包含于 $G$，
故它们的长度和不超过 $m(G)$；上式中的两个估计因此都已显式核对。
Darboux 判据遂给出 $f$ Riemann 可积。
\end{proof}

\begin{example}[Thomae 函数与 Dirichlet 函数]
在 $[0,1]$ 上令
\[
 t(x)=\begin{cases}q^{-1},&x=p/q\text{ 为既约分数},\\0,&x\notin\Q,
 \end{cases}
 \qquad d(x)=\1_{\Q}(x).
\]
$t$ 在每个有理点不连续，因为无理数列趋近该点而函数值恒为零。若 $x$ 无理，给定
$\varepsilon>0$，分母不超过 $1/\varepsilon$ 的既约分数只有有限多个；取一个避开这些点的
$x$ 的小邻域，其中 $t<\varepsilon$，故 $t$ 在 $x$ 连续。因此 $t$ 的不连续点集可数，
由 \cref{theorem:4-1-2-3} 它 Riemann 可积；每个非退化区间含无理数，所以所有 Darboux 下和为
$0$，其积分也为 $0$。函数 $d$ 在每个区间内同时取 $0$ 与 $1$，故每个分割的上、下和分别为
$1$ 和 $0$；它处处不连续，其不连续集测度为 $1$，因而不 Riemann 可积。这一对照也说明
“坏点稠密”并非判据。
\end{example}

球平均的形状可以换成上述中心立方体，也可以在有适当覆盖性质的微分基中推广；任意一列缩向点的
可测集合则不受本定理保证。下一小节把被平均的函数换成测度，检验局部质量比究竟恢复哪一部分信息。

\begin{exercise}
直接由定义证明连续点是 Lebesgue 点，并证明 Lebesgue 点处 $A_rf(x)\to f(x)$。
\end{exercise}
\begin{exercise}
计算 $E=[0,1]\cup[2,3]$ 的每个点的密度（只需讨论区间内部、端点和外部）。
\end{exercise}
\begin{exercise}
在 Riemann 可积判据的必要性证明中，详细证明若 $x\in D_k$ 位于分割小区间内部，则该区间
振荡至少为 $1/k$。
\end{exercise}
\begin{exercise}
对 $f=\1_E$ 应用 Lebesgue 微分定理，完整推出 Lebesgue 密度定理；须分别说明 $x\in E$ 与
$x\notin E$ 时平均值的极限。
\end{exercise}
\begin{exercise}
独立完成 Riemann 可积判据的充分性方向：把不连续点放进小测度开集，并在其补集上用有限连续性
邻域控制 Darboux 上、下和之差。
\end{exercise}

\subsection{测度的局部密度与抽象微分}\label{subsec:4-1-3}
\index{测度的密度}\index{Lebesgue 分解}

函数导数是增量之比的极限。若把函数换成测度，最直接的候选便是
\[
 \frac{\nu(B(x,r))}{m(B(x,r))}.
\]
三类测度说明了这一比值的不同表现。若 $\nu=fm$ 且 $f$ 连续，该比值是 $f$ 的球平均，因而趋于
$f(x)$；若 $\nu=\delta_0$，它在 $0$ 发散，在其他点最终为零；若质量集中在 Cantor 集上，则它在
Lebesgue 几乎处处为零。由此可见，局部体积比恢复绝对连续部分，却不能以普通函数记录奇异质量。

\begin{definition}[测度密度比]\label{def:4-1-3-1}
对 $\R^n$ 上的局部有限正 Borel 测度 $\nu$，定义
\[
 D_r\nu(x)=\frac{\nu(B(x,r))}{m(B(x,r))}.
\]
若 $\lim_{r\downarrow0}D_r\nu(x)$ 存在，则称它为 $\nu$ 相对于 Lebesgue 测度的对称密度。
对局部有限有符号或复测度，同一定义用于其有限取值的球质量。
\end{definition}

\begin{definition}[中心微分基与相对微分]\label{def:4-1-3-2}
设 $\mu$ 是局部有限正 Radon 测度。一个中心微分基 $\mathcal D$ 为每个 $x$ 指定一族含 $x$ 的有界
Borel 集 $\mathcal D(x)$，且其中集合的直径可以任意小。记 $V\to x$ 表示
$V\in\mathcal D(x)$、$\mu(V)>0$ 且 $\diam V\to0$。

若 $\nu$ 是局部有限有符号 Radon 测度且 $\nu\ll\mu$，则沿该微分基考察相对比值
\[
 \frac{\nu(V)}{\mu(V)}.
\]
其极限称为 $\nu$ 关于 $\mu$ 沿 $\mathcal D$ 的相对微分。分母所用的测度、集合缩向点的方式以及
“几乎处处”所相对的测度都是定义的一部分。
\end{definition}

\begin{definition}[上、下密度]\label{def:4-1-3-3}
对正测度 $\nu$，定义
\[
 \overline D\nu(x)=\limsup_{r\downarrow0}D_r\nu(x),\qquad
 \underline D\nu(x)=\liminf_{r\downarrow0}D_r\nu(x).
\]
二者允许取 $+\infty$。
\end{definition}

\begin{example}[连续密度]\label{ex:4-1-3-1}
若 $\nu=fm$，其中 $f\ge0$ 连续，则
\[
 D_r\nu(x)=A_rf(x)\longrightarrow f(x)
\]
在每个点成立。若只假设 $f\in L^1_{\mathrm{loc}}$，同一计算仍成立，只是极限由
\cref{theorem:4-1-2-1} 保证在几乎每一点成立。
\end{example}

\begin{example}[Dirac 测度]\label{ex:4-1-3-2}
对 $x\ne0$，当 $0<r<|x|$ 时 $\delta_0(B(x,r))=0$，故
$D_r\delta_0(x)$ 最终为零。在 $x=0$，
\[
 D_r\delta_0(0)=\frac1{\omega_nr^n}\longrightarrow+\infty.
\]
因此 $D_r\delta_0\to0$ 对 $m$-几乎每一点成立，而 $\delta_0(\R^n)=1$。密度为零几乎处处
不能推出原测度为零。
\end{example}

\begin{example}[Cantor 测度]\label{ex:4-1-3-3}
令 $C\subset[0,1]$ 为三分 Cantor 集，$\mu_C$ 为标准 Cantor 概率：每个第 $k$ 级基本区间
质量为 $2^{-k}$。记 $\alpha=\log2/\log3$。若
$3^{-(k+1)}<r\le3^{-k}$ 且 $x\in C$，含 $x$ 的第 $k+1$ 级基本区间长度小于 $r$，故整个
区间包含在 $B(x,r)$ 中，故
\[
 \mu_C(B(x,r))\ge2^{-(k+1)}\ge\frac12r^\alpha.
\]
另一方面，两个不同的第 $k$ 级基本区间之间的距离至少为 $3^{-k}$，而
$B(x,r)$ 的长度不超过 $2\cdot3^{-k}$，所以它至多与三个第 $k$ 级基本区间相交。从而
\[
 \mu_C(B(x,r))\le3\cdot2^{-k}<6r^\alpha,
\]
最后一步用了 $r>3^{-(k+1)}$。因此可以显式取 $c=1/2$ 和 $C=6$，使
\[
 c r^\alpha\le\mu_C(B(x,r))\le C r^\alpha
 \qquad(x\in C,\ 0<r<1).
\]
这里使用了 $2^{-k}=(3^{-k})^\alpha$。于是对 $x\in C$，
\[
 \frac14r^{\alpha-1}
 \le\frac{\mu_C(B(x,r))}{m(B(x,r))}
 \le3r^{\alpha-1},
 \qquad
 \frac{\mu_C(B(x,r))}{m(B(x,r))}\longrightarrow+\infty.
\]
可是 $m(C)=0$，且 $C$ 闭；对 $m$-几乎每个 $x$，事实上对每个 $x\notin C$，充分小的球不碰
$C$，所以该比值最终为零。同一测度在 $m$-几乎处处不可见，却在自身支撑上的每一点高度集中。
\end{example}

\begin{figure}[htbp]
\centering
\begin{tikzpicture}[x=1.05cm,y=.82cm]
  \draw[->,BookInk!75] (0,0)--(8.1,0)
    node[right,book label] {$s=-\log r$（向右即 $r\downarrow0$）};
  \draw[->,BookInk!75] (0,-.25)--(0,5.2)
    node[above,book label] {$\log D_r$};
  \draw[BookAccent,very thick] (.5,1.25)--(7.4,1.25)
    node[right,book label] {$fm$：$\log f(x)$};
  \draw[BookWarning,very thick] (.5,.35)--(7.3,4.7)
    node[right,book label] {$\delta_0$ 在 $0$：斜率 $n$};
  \draw[BookWarm,very thick,dashed] (.5,.7)--(7.3,3.2)
    node[right,book label] {$\mu_C$ 在 $C$：斜率 $1-\alpha$};
  \node[book node,align=left] at (3.8,4.55)
    {连续密度趋于有限值；\\奇异质量在自身支撑上发散};
\end{tikzpicture}
\caption{三种局部质量比置于同一对数尺度。水平线按 $f(x)>0$ 绘制；
在 Dirac 或 Cantor 支撑之外，相应比值对充分小的 $r$ 恒为零。图展示渐近阶，具体结论仍由
例~\ref{ex:4-1-3-1} 至例~\ref{ex:4-1-3-3} 的计算给出。}
\label{fig:4-1-density-scale}
\end{figure}

\begin{theorem}[绝对连续测度微分]\label{theorem:4-1-3-1}
若 $f\in L^1_{\mathrm{loc}}(\R^n)$ 且 $\nu=fm$，则
\[
 \lim_{r\downarrow0}\frac{\nu(B(x,r))}{m(B(x,r))}=f(x)
\]
对 $m$-几乎每个 $x$ 成立。
\end{theorem}

\begin{proof}
对每个球，
\[
 \frac{\nu(B(x,r))}{m(B(x,r))}
 =\frac1{m(B(x,r))}\int_{B(x,r)}f(y)\,\dd y=A_rf(x).
\]
结论正是 \cref{theorem:4-1-2-1}。
\end{proof}

为了处理奇异部分，先记录弱型证明并不需要被平均的质量具有密度。若 $\sigma$ 是有限正 Radon 测度，
定义
\[
 M\sigma(x)=\sup_{r>0}\frac{\sigma(B(x,r))}{m(B(x,r))}.
\]
固定 $r$ 时 $x\mapsto\sigma(B(x,r))$ 下半连续：若 $x_j\to x$，则对每个 $0<s<r$，
$B(x,s)\subset B(x_j,r)$ 当 $j$ 充分大，故
\[
 \liminf_j\sigma(B(x_j,r))\ge\sigma(B(x,s));
\]
再令 $s\uparrow r$，便知每个固定半径的质量函数下半连续；取上确界后
$M\sigma$ 仍下半连续，因而可测且其严格上水平集为开集。为把函数情形的弱型估计移到这里，记
$E_\lambda=\{M\sigma>\lambda\}$，并固定紧集 $K\subset E_\lambda$。对每个 $x\in K$ 选取
见证球 $B(x,r_x)$，使
\[
  \sigma(B(x,r_x))>\lambda m(B(x,r_x)).
\]
由紧性可先取有限子覆盖，再对这个有限球族应用
\cref{lemma:4-1-1-1}，得到两两不交的见证球
$B_1,\ldots,B_N$，且 $K\subset\bigcup_{j=1}^N5B_j$。于是
\[
 \begin{aligned}
 m(K)&\le\sum_{j=1}^Nm(5B_j)
      =5^n\sum_{j=1}^Nm(B_j)\\
 &<\frac{5^n}{\lambda}\sum_{j=1}^N\sigma(B_j)
 \le\frac{5^n}{\lambda}\sigma(\R^n).
 \end{aligned}
\]
最后对 $E_\lambda$ 的紧子集取上确界；$E_\lambda$ 是开集，Lebesgue 测度的内正则性给出
\begin{equation}\label{eq:4-1-measure-maximal}
 m\{M\sigma>\lambda\}\le\frac{5^n}{\lambda}\sigma(\R^n).
\end{equation}
需要注意，若奇异测度集中在零测集 $N$ 上，取开集 $U\supset N$ 并不能使
$\nu_s(U)$ 变小，因为 $U$ 仍含有全部奇异质量。下面的证明将改为截去承载绝大部分质量的紧集，
再对剩余测度应用极大估计。

\begin{theorem}[Lebesgue 分解的局部图像]\label{theorem:4-1-3-2}
设 $\nu$ 是 $\R^n$ 上的局部有限正 Radon 测度，且
\[
 \nu=fm+\nu_s,\qquad \nu_s\perp m,\qquad f\in L^1_{\mathrm{loc}}(\R^n)
\]
是它的 Lebesgue 分解。则
\[
 D_r\nu(x)\longrightarrow f(x)
\]
对 $m$-几乎每个 $x$ 成立。特别地，$D_r\nu_s(x)\to0$ 对 $m$-几乎每个 $x$ 成立。
\end{theorem}

\begin{proof}
绝对连续部分由 \cref{theorem:4-1-3-1} 处理，只需证明关于 $\nu_s$ 的最后一句。
取 Borel 零测集 $N$，使 $\nu_s(\R^n\setminus N)=0$。固定正整数 $R$，令
$W=B(0,R+1)$ 及 $\sigma=\nu_s|_W$。这是有限正 Radon 测度。给定 $\varepsilon>0$，
内正则性给出紧集 $K\subset N\cap W$，使
\[
 \sigma(W\setminus K)<\varepsilon.
\]
令 $\tau=\sigma|_{W\setminus K}$，则 $\tau(\R^n)<\varepsilon$。

若 $x\in B(0,R)\setminus K$，由于 $K$ 闭且 $B(0,R)$ 到 $W^c$ 的距离为 $1$，存在
$r_x>0$，使 $0<r<r_x$ 时 $B(x,r)\subset W\setminus K$。若 $K\ne\varnothing$，可取
\[
 r_x=\frac12\min\{1,\dist(x,K)\}>0.
\]
若 $K=\varnothing$，则直接取 $r_x=1/2$。
因此
\[
 \nu_s(B(x,r))=\sigma(B(x,r))=\tau(B(x,r))
 \qquad(0<r<r_x).
\]
对 $\lambda>0$，这说明
\[
 \left\{x\in B(0,R)\setminus K:
   \limsup_{r\downarrow0}D_r\nu_s(x)>\lambda\right\}
 \subset\{M\tau>\lambda\}.
\]
由于 $m(K)=0$，由 \eqref{eq:4-1-measure-maximal} 得
\[
 m\left\{x\in B(0,R):
   \limsup_{r\downarrow0}D_r\nu_s(x)>\lambda\right\}
 \le\frac{5^n\varepsilon}{\lambda}.
\]
令 $\varepsilon\downarrow0$，左侧集合为零测集。再对 $\lambda=1/k$ 取可数并，得到
$\limsup_{r\downarrow0}D_r\nu_s=0$ 几乎处处于 $B(0,R)$。密度比非负，所以其下极限不小于
$0$，从而极限为 $0$。最后令 $R\uparrow\infty$，得到全空间结论。

在绝对连续部分与奇异部分各自的满测集合之交上，
$D_r\nu=D_r(fm)+D_r\nu_s\to f+0$，证明完成。
\end{proof}

\begin{proposition}[相对密度与 Radon--Nikodym 导数]\label{proposition:4-1-3-3}
设 $\mathcal D$ 是关于局部有限正 Radon 测度 $\mu$ 的一个中心微分基，并假定它能够微分
$L^1_{\mathrm{loc}}(\mu)$：即对每个 $h\in L^1_{\mathrm{loc}}(\mu)$，沿 $V\in\mathcal D(x)$
缩向 $x$ 时
\[
 \frac1{\mu(V)}\int_Vh\,\dd\mu\longrightarrow h(x)
\]
对 $\mu$-几乎每个 $x$ 成立。若 $\nu=h\mu\ll\mu$，则沿同一微分基
\[
 \frac{\nu(V)}{\mu(V)}\longrightarrow h(x)
 =\frac{\dd\nu}{\dd\mu}(x)
\]
对 $\mu$-几乎每个 $x$ 成立。
\end{proposition}

\begin{proof}
在 $\mu(V)>0$ 的基集合上，Radon--Nikodym 表示直接给出
\[
 \frac{\nu(V)}{\mu(V)}
 =\frac1{\mu(V)}\int_Vh\,\dd\mu.
\]
右端由题设的微分基性质趋于 $h(x)$。因此，除 Radon--Nikodym 表示外，还必须假设
$\mathcal D$ 能够微分 $L^1_{\mathrm{loc}}(\mu)$。
\end{proof}

\begin{remark}[Besicovitch 微分定理（选读）]\label{proposition:4-1-3-4}
在 $\R^n$ 中，对任意局部有限正 Radon 测度 $\mu$，中心球基能够微分
$L^1_{\mathrm{loc}}(\mu)$ 中的函数。因此 \cref{proposition:4-1-3-3} 的相对密度结论对中心球成立。
\end{remark}

其证明需要 Besicovitch 覆盖定理：带指定中心的球族可分成至多 $N(n)$ 个不交子族，并仍覆盖中心集。
该覆盖结论产生相对于 $\mu$ 的极大弱型估计，再与连续逼近结合即可完成证明。这里不展开该覆盖定理。
Lebesgue 测度版的 $5r$ 论证不能直接用于一般 $\mu$，因为五倍球的 $\mu$-质量与原球质量之间通常
没有统一比值。本章后续结论不使用这条选读定理；它只记录一般相对微分所需的额外覆盖输入。

\begin{remark}[混合测度的局部密度]
取 $0\le p\le1$，令
\[
 \nu=p\delta_0+(1-p)\1_{(0,1)}m.
\]
它的总质量为 $1$。对 $m$-几乎每个 $x$，局部质量比只恢复
$(1-p)\1_{(0,1)}(x)$；原子质量 $p$ 在 $x=0$ 的比值中爆炸，却在 Lebesgue 几乎处处消失。
因此，总质量为一并不蕴含测度具有概率密度；后者要求该测度绝对连续于 $m$。
\end{remark}

局部质量比既可解释概率密度，也可用来研究数值密度估计；不过有限样本产生的是带带宽的近似比值，
不能把本节的几乎处处极限直接当作有限样本误差界。另一方面，一维有界变差函数的导数测度会同时含有
绝对连续和奇异部分，下一节将给出相应分解。

\begin{figure}[htbp]
\centering
\begin{tikzpicture}[node distance=13mm and 20mm]
  \node[book strong node] (nu) {$\nu$};
  \node[book node,below left=of nu] (ac) {铺在体积上的部分\\$fm$};
  \node[book warning node,below right=of nu] (sing) {零测骨架上的部分\\$\nu_s$};
  \node[book warm node,below=of ac] (ratio) {一般点的小球比值\\$\longrightarrow f(x)$};
  \node[book node,below=of sing] (zero) {$m$-a.e. 比值\\$\longrightarrow0$};
  \draw[book arrow] (nu)--(ac);
  \draw[book arrow] (nu)--(sing);
  \draw[book accent arrow] (ac)--(ratio);
  \draw[book dashed arrow] (sing)--(zero);
\end{tikzpicture}
\caption{Lebesgue 分解的局部图像。小球体积比恢复 $fm$ 的密度；奇异质量并未消失，
只是对 Lebesgue 几乎每个中心不可见。}
\label{fig:4-1-local-density-decomposition}
\end{figure}

\begin{exercise}
证明对 $x\ne0$，$D_r\delta_0(x)$ 最终为零；计算在 $x=0$ 的发散阶。
\end{exercise}
\begin{exercise}
把 \eqref{eq:4-1-measure-maximal} 的证明完整写出，并指出只在何处使用 $\sigma$ 的有限性。
\end{exercise}
\begin{exercise}
设 $S^{n-1}$ 上的表面测度作为 $\R^n$ 中的测度。证明它相对于 $n$ 维 Lebesgue 测度的密度
在 $m$-几乎处处为零，并说明球面上的中心点为何不与此矛盾。
\end{exercise}
\begin{exercise}
若 $\nu=fm$、$f\in L^1_{\mathrm{loc}}(\R^n)$，把密度比逐行化成球平均，并由 Lebesgue
微分定理推出 $D_r\nu(x)\to f(x)$ 对 $m$-几乎每个 $x$ 成立。
\end{exercise}
\begin{exercise}
给出一个非零奇异测度，使 $D_r\nu(x)\to0$ 对 $m$-几乎每点成立；分别用其总质量与局部比值解释
为何“密度几乎处处为零”不推出测度为零。
\end{exercise}

局部质量比已经把 Radon--Nikodym 密度变成可观察的极限，同时也暴露了奇异部分。实线上的奇异质量
常常可以编码进一个单调或有界变差函数；下一节将从分割上的总变差开始，把绝对连续、跳跃与奇异连续
三种行为逐一拆开。
